\begin{table}[!ht]
	\caption{Policy Implications of Documented Exaggeration Risk}
	\label{tab:policy_implications}
	\centering
		\begin{threeparttable}
                		\begin{tabular}{lccccc}
                    		\toprule
                    			 & Exaggeration & Expected  & Adjusted  & Marginal\\ 
			 		 & ratio & Marginal & Marginal  & Benefits \\
			 		 &  & Benefits & Benefits & to Costs \\ 
			 		 &  &  (B\$/year) & (B\$/year) & Ratio \\ 
				\midrule
Epidemiology median & 1.3 & 13.5 to 29 & 10.4 to 22.3 & 27 to 57\\
Epidemiology 25th pctile & 1.9 & 13.5 to 29 & 7.1 to 15.3 & 18 to 39\\
Causal inference median & 1.7 & 13.5 to 29 & 7.9 to 17.1 & 20 to 44\\
Causal inference 25th pctile & 2.0 & 13.5 to 29 & 6.8 to 14.5 & 17 to 37\\
IV vs OLS median & 4.5 & 13.5 to 29 & 3 to 6.4 & 8 to 17\\
Causal vs standard example & 5.4 & 13.5 to 29 & 2.5 to 5.4 & 6 to 14\\
                    		\bottomrule
                \end{tabular}
                \begin{tablenotes}
                    \footnotesize
                    \item \textit{Notes}: Adjusted marginal benefits represent the benefits under the hypothesis that health effects were exaggerated: due to the linearity of the relationship between health effects and monetized benefits, they are simply computed as the expected benefits of tightening the policy from 10 to 9 $\mu\text{g/m}^3$ divided by the exaggeration ratios found in our analysis of the literature. The various exaggeration ratios considered are drawn from the analysis of the literature above, except for the last one which comes from the comparison between \cite{schwartz_estimating_2015} and \cite{zanobetti_effect_2009}.
                \end{tablenotes}
            \end{threeparttable}
\end{table}